\section{Unified Sovereign Intelligence Contract}
\label{sec:unified-sovereign-intelligence-contract}

This section records completion of Sovereign Business Brain roadmap tasks
\texttt{SBB-0301}--\texttt{SBB-0308}.  AION now has one provider-neutral routing
boundary for deterministic capabilities, retrieval, evidence services, local and
remote models, and agent harnesses.  A model remains replaceable compute: it does
not become the customer brain, own memory, grant permission or execute a business
action by itself.

\subsection{Portable contract family}

The contract family is published as portable JSON Schema and as typed runtime
models.  It contains four primary envelopes:

\begin{description}
  \item[Intelligence Request] identifies the customer brain, idempotent request,
  task class, permitted input, required capabilities, output schema, route budget,
  disclosure policy and final delivery target.
  \item[Intelligence Response] returns a structured result, chosen adapter,
  normalized attempts, delivery status and Intelligence Route Receipt.
  \item[Stream Event] provides ordered accepted, attempt, fallback, delivery and
  completion events without copying the original prompt into routing telemetry.
  \item[Adapter Manifest] describes an adapter's kind, location, provider, model,
  capabilities, destination, residency, supported data classification and expected
  quality, cost, latency and energy.
\end{description}

All five adapter kinds---\texttt{capability}, \texttt{model}, \texttt{retrieval},
\texttt{evidence} and \texttt{agent\_harness}---use this same boundary.  Consequently,
an open model on the customer laptop, a model in the customer's private cloud, a
public evidence service and an optional premium API can be substituted without
moving the canonical business brain.

\subsection{Governed routing sequence}

\begin{verbatim}
Pilot / Boardroom request
        |
        v
AION identity + task + schema validation
        |
        v
Context-size and route-budget admission
        |
        v
Candidate capability check
        |
        +-- local adapter -----------------------------+
        |                                              |
        +-- customer cloud / external API              |
              disclosure + destination + residency    |
              field minimisation                      |
                                                       v
                                      structured output validation
                                                       |
                                      quality / usage verification
                                                       |
                                      Pilot presentation OR
                                      governed capability authority
                                                       |
                                      hash-only route receipt
\end{verbatim}

The caller supplies an ordered candidate list.  The router evaluates it
deterministically and cannot add an undeclared provider.  A fallback therefore
cannot escape the customer's privacy, residency or cost policy simply because a
preferred model failed.

\subsection{External disclosure and context minimisation}

Local adapters may receive the admitted local request.  A customer-cloud or external
adapter is rejected unless external disclosure is enabled and its exact destination
and residency occur in the request's allowlists.  The request must also name the
fields that may leave the brain.  A separate request copy is constructed for the
external adapter containing only those named top-level fields; all other local input
is removed before the adapter callable is invoked.

The manifest additionally declares the highest data classification the destination
may receive.  Public, internal, confidential and restricted data are ordered, and a
route is ineligible when the request classification exceeds the adapter's declared
limit.  These checks happen before network-capable code receives the request.

\subsection{Multi-dimensional budgets}

Each request carries maximum context bytes, financial cost, latency, energy and
attempt count together with a minimum acceptable quality.  Estimated adapter use is
checked before an attempt.  Reported use and structured-output quality are checked
afterward.  A cheap but inadequate result cannot satisfy a high-quality request, and
a powerful provider cannot be selected when its cost, location or disclosure exceeds
the owner's boundary.

The budget is a route authority rather than a reporting hint.  Context overflow fails
before execution.  Ineligible candidates are recorded as skipped with a machine-readable
reason.  Attempts stop at the declared ceiling.

\subsection{Health, fallback and idempotency}

Adapter success resets its failure counter.  Repeated failures open a persisted
circuit for a bounded cooling period, preventing a broken provider from being called
for every subsequent request.  The router then evaluates the next candidate under
the same original disclosure and budget rules.

The request identifier is an idempotency key.  Replaying the identical request returns
the recorded response without calling the adapter or delivery sink again.  Reusing the
identifier with changed content fails closed.  This protects both paid inference and
consequential downstream capabilities from accidental duplicate execution.

\subsection{Receipts without prompt retention}

Every terminal outcome emits an Intelligence Route Receipt containing request and
result hashes, normalized route attempts, disclosure policy and budget.  It explicitly
records that raw prompts and raw results are not retained in the receipt.  When a
Sovereign Brain boundary is supplied, the receipt is placed in its canonical receipt
store under a content-derived identifier.  This makes routing auditable without
turning telemetry into a second store of private conversations.

\subsection{Presentation and action separation}

A successful result has two possible destinations.  A Pilot presentation sink may
format the structured output for an authorized phone, television, desktop or browser.
A consequential result must instead enter a governed capability sink.  That path
requires both an exact capability identifier and an approval receipt; absence of
either authority fails delivery.  The router passes a proposal to the capability
authority but cannot manufacture permission, bypass the capability policy or claim
that an external action succeeded.

This preserves the core sequence:

\begin{center}
\textbf{understand $\rightarrow$ route $\rightarrow$ validate $\rightarrow$ approve
$\rightarrow$ execute $\rightarrow$ verify $\rightarrow$ learn.}
\end{center}

\subsection{Verification and release record}

Automated verification covers schema publication, every adapter kind and local Pilot
delivery.  It also covers private cloud disclosure, destination and residency controls,
removal of undeclared fields, resource budgets, output checking and persistent circuits.
Further tests prove fallback containment, idempotent replay, changed request rejection,
approved capability delivery, ordered streaming events and receipt persistence without
raw content.

The implementation is packaged in AION Sovereign Brain Bootstrap~0.6.0 and Pilot
Fabric~0.58.0.  It does not silently migrate every historical direct model call; those
call sites can now be moved behind this canonical contract incrementally.  The next
roadmap phase is the signed model catalogue and qualification system, which decides
which concrete models are eligible to register as adapters.
